\chapter{Conclusion}

This project set out to develop \textit{Tourlingo}, a web-based assistant that helps tourists overcome language and cultural barriers when exploring new places. Within three months, the team delivered a functional prototype that combined mapping, translation, and cultural guidance into one system.

\section{Critical Evaluation}
Overall, the project achieved the Minimum Viable Product (MVP) defined in the PRD. The OpenStreetMap based interface, enhanced with Overpass queries, provided with nearby places for essential services and recommendations. Custom pin icons and side panels improved clarity and usability. Integration with LibreTranslate enabled dynamic translations of directions and cultural tips, although occasional inaccuracies showed the limits of relying on free, open-source APIs.  

The backend evolved iteratively, with horizontal scaling and database caching reducing response times to near real-time. On the frontend, Next.js and TailwindCSS supported a responsive and modular design that was easy to adapt during user testing. Two departmental testathons were particularly valuable, as feedback from over thirty participants led to concrete fixes such as smoother panel navigation and improved search responsiveness. With hindsight, richer POI data (e.g. images, opening hours) could have enhanced the experience, but this was restricted by the APIs chosen.

\section{Summary}
In summary, Tourlingo achieved its core objectives and demonstrated the value of combining microservice architecture, modern web technologies, and user-centred evaluation. The system is not only a working prototype for the client but also evidence that a small team can design and deploy a scalable travel assistant within strict constraints. While further work is needed to enrich data and expand features, the project provides a strong foundation for future development.
